pdfoutput=1
\pdfoutput=1
\section{Theorem 5: Cluster Consistency of State Discovery}
\label{sec:S5}

This section proves Theorem~5, which establishes that under the
\textbf{strong features} regime (large separation $\Delta_{\min}$ relative
to noise), Lloyd's $k$-means with $O(\log n)$ random restarts recovers
the true state partition with probability tending to one exponentially fast.
Theorem~5 is the positive counterpart to the negative result of Theorem~2
(weak features $\to$ failure) and Theorem~4' (minimax lower bound), and
it complements the F1 bound from Theorem~1 \eqref{eq:main:F1} and the
weak-feature F1 bound from Theorem~2 \eqref{eq:main:weakF1}.
Together these results complete the phase diagram: strong features imply
state discovery succeeds; weak features imply SCX cannot exceed the loss
baseline; and the threshold is minimax optimal as established in
Theorem~4' \eqref{eq:main:minimax}.

% ---------------------------------------------------------------------------
\subsection{Setup: Fixed-$K$ Generative Model with Strong Features}
\label{sec:S5:setup}
% ---------------------------------------------------------------------------

We formalize the data-generating process and the clustering objective.

\begin{assumption}[Generative model]
\label{ass:S5:gen}
Let $(\Omega,\cF,\Pp)$ be a probability space. We observe $n$ i.i.d.\ copies
$\phi(x_i)\in\R^{d_\phi}$ of the feature vector, generated as
\begin{equation}
\phi(x) = \mu_{s(x)} + \varepsilon,
\label{eq:S5:gen}
\end{equation}
where
\begin{itemize}
\item $s(x)\in\{1,\dots,K\}$ is the unobserved true state,
\item $K$ is \textbf{fixed} (does not grow with $n$),
\item $\muk\in\R^{d_\phi}$ is the feature mean (center) for state $k$,
\item $\varepsilon\in\R^{d_\phi}$ is a zero-mean sub-Gaussian random vector
      with variance proxy $\sigma^2>0$, independent of $s(x)$:
      \[
      \E[\varepsilon]=0,\qquad
      \E\bigl[\exp(t\cdot u^\top\varepsilon)\bigr]\le\exp(\sigma^2 t^2/2),\quad
      \forall u\in\mathbb{S}^{d_\phi-1},\;\forall t\in\R.
      \]
\end{itemize}
\end{assumption}

\begin{assumption}[Cluster structure]
\label{ass:S5:cluster}
The true partition is $\cstar=\{C_1^*,\dots,C_K^*\}$ with
$C_k^*=\{x\in\cX: s(x)=k\}$.  The true centers are
$\muk=\E[\phi(X)\mid S=k]$.  The population proportions are
$\pi_k=P(S=k)>0$ for all $k$, and we denote
$\pi_{\min}=\min_k\pi_k$.  The empirical per-state sample sizes are
$n_k=\sum_{i=1}^n\ind\{s(x_i)=k\}$ with $n_{\min}=\min_k n_k$.
\end{assumption}

\begin{assumption}[Strong separation (fixed $\Delta_{\min}$)]
\label{ass:S5:sep}
The minimum distance between distinct cluster centers is a fixed positive
constant (it does \emph{not} scale with $n$ or $K$):
\[
\Delta_{\min} \;=\; \min_{i\neq j}\norm{\mu_i-\mu_j}_2 \;>\;0.
\]
We assume the \textbf{strong separation regime}
\begin{equation}
\frac{\Delta_{\min}^2}{\sigma^2 d_\phi}\;\ge\;C_0,
\label{eq:S5:strongsep}
\end{equation}
where $C_0$ is a universal constant large enough to guarantee
$\norm{\theta^*-\mu}<\Delta_{\min}/8$ in Lemma~\ref{lem:S5:pop} below.
This is the ``strong features'' regime: the signal-to-noise ratio is
sufficiently high that cluster centers are distinguishable despite noise
and feature dimensionality.
\end{assumption}

\begin{definition}[$k$-means objective]
\label{def:S5:kmeans}
For a set of $K$ candidate centers
$\theta=\{\theta_1,\dots,\theta_K\}\subset\R^{d_\phi}$, define the
\textbf{empirical $k$-means risk}
\[
W_n(\theta)=\frac1n\sum_{i=1}^n\min_{k\in[K]}\norm{\phi(x_i)-\theta_k}_2^2,
\]
and the \textbf{population $k$-means risk}
\[
W(\theta)=\E\Bigl[\min_{k\in[K]}\norm{\phi(X)-\theta_k}_2^2\Bigr].
\]
Let
$\theta^*\in\argmin_{\theta:|\theta|=K}W(\theta)$ be the population
minimizer and $\htheta_n\in\argmin_{\theta:|\theta|=K}W_n(\theta)$ the
global empirical minimizer, both defined up to label permutation.
When we write $\norm{\htheta_n-\theta^*}$, we mean the distance after
optimal permutation matching.
\end{definition}

\begin{definition}[Norms on center sets]
For a $K$-center set $\theta$, let
\[
\norm{\theta}_\infty = \max_{k\in[K]}\norm{\theta_k}_2,
\qquad
\norm{\theta-\theta'} = \max_{k}\norm{\theta_k-\theta'_k}_2.
\]
\end{definition}

% ---------------------------------------------------------------------------
\subsection{Lemma~1: Population Minimizer Proximity}
\label{sec:S5:lem1}
% ---------------------------------------------------------------------------

\begin{lemma}[Population minimizer proximity]
\label{lem:S5:pop}
Under Assumptions~\ref{ass:S5:gen}--\ref{ass:S5:sep}, the population
$k$-means objective $W(\theta)$ has a minimizer $\theta^*$ (unique up
to label permutation).  For any permutation $\pi$ that optimally matches
$\theta^*$ to the true centers $\mu$,
\[
\norm{\theta^*_{\pi(k)}-\muk}_2 \;\le\; \epsilon_{\mathrm{pop}}
\;<\; \frac{\Delta_{\min}}{8},
\]
where
\[
\epsilon_{\mathrm{pop}} \;=\;
\frac{C_2\,K\,(M+\sigma\sqrt{d_\phi})}{\pi_{\min}}
\cdot \exp\!\left(-\frac{\Delta_{\min}^2}{8\sigma^2}\right),
\]
and $C_2$ is an absolute constant, $M=\max_k\norm{\muk}_2$.
Under \eqref{eq:S5:strongsep} with $C_0$ sufficiently large,
$\epsilon_{\mathrm{pop}}<\Delta_{\min}/8$.
\end{lemma}

\begin{proof}[Proof of Lemma~\ref{lem:S5:pop}]
The proof proceeds in five steps.

\noindent\textbf{Step 1 (Existence).}
$W(\theta)$ is continuous (composition of continuous functions) and
coercive ($W(\theta)\to\infty$ as $\norm{\theta}_\infty\to\infty$,
since the quadratic penalty dominates), hence attains a minimum on any
compact set containing a sufficiently large ball.  By strict convexity
of $\theta_k\mapsto\E[\norm{\phi-\theta_k}^2\mid\phi\in V_k(\theta)]$
on each Voronoi cell (defined below), any two minimizers must be related
by a label permutation.  Thus $\theta^*$ is unique up to permutation.

\noindent\textbf{Step 2 (Self-consistency equation).}
For a set of centers $\theta$, define the Voronoi cell of center $j$:
\[
V_j(\theta)=\bigl\{\phi\in\R^{d_\phi}:
\norm{\phi-\theta_j}_2\le\min_{\ell\neq j}\norm{\phi-\theta_\ell}_2\bigr\}.
\]
At points where Voronoi cells have positive probability and are
well-defined, the gradient of $W$ with respect to $\theta_j$ exists:
\[
\frac{\partial}{\partial\theta_j}W(\theta)
= -2\,\E\bigl[(\phi-\theta_j)\cdot\ind\{\phi\in V_j(\theta)\}\bigr].
\]
Setting this to zero at the minimizer $\theta^*$ gives the
\textbf{self-consistency condition}: for each $j$,
\[
\theta^*_j = \frac{\E[\phi\cdot\ind\{\phi\in V_j(\theta^*)\}]}
                 {P(\phi\in V_j(\theta^*))}
           = \E[\phi\mid\phi\in V_j(\theta^*)]. \tag{1}
\]

\noindent\textbf{Step 3 (True centers approximately satisfy self-consistency).}
Consider the Voronoi cells induced by the true centers $\mu=\mutrue$.
For each $k$, define the one-step candidate
\[
\ttheta_k = \E[\phi\mid\phi\in V_k(\mu)].
\]
We bound $\norm{\ttheta_k-\muk}$.

\noindent\textit{Claim (mis-assignment probability bound).}
For $j\neq k$,
\[
P\bigl(\mu_j+\varepsilon\in V_k(\mu)\mid S=j\bigr)
\;\le\; \exp\!\left(-\frac{\norm{\mu_j-\muk}^2}{8\sigma^2}\right)
\;\le\; \exp\!\left(-\frac{\Delta_{\min}^2}{8\sigma^2}\right).
\]

\noindent\textit{Proof of claim.}
The event $\mu_j+\varepsilon\in V_k(\mu)$ requires
$\norm{\mu_j+\varepsilon-\muk}\le\norm{\mu_j+\varepsilon-\mu_j}$, i.e.\
$\norm{(\mu_j-\muk)+\varepsilon}\le\norm{\varepsilon}$.  Squaring:
\[
\norm{\mu_j-\muk}^2 + 2\varepsilon^\top(\mu_j-\muk) + \norm{\varepsilon}^2
\;\le\; \norm{\varepsilon}^2,
\]
hence $2\varepsilon^\top(\mu_j-\muk) \le -\norm{\mu_j-\muk}^2$.
Since $v=(\mu_j-\muk)/\norm{\mu_j-\muk}$ is a unit vector,
$\varepsilon^\top v$ is sub-Gaussian with parameter $\sigma^2$, so
$2\varepsilon^\top(\mu_j-\muk)$ is sub-Gaussian with parameter
$4\sigma^2\norm{\mu_j-\muk}^2$.  By the sub-Gaussian tail bound,
\[
P\bigl(2\varepsilon^\top(\mu_j-\muk)\le -\norm{\mu_j-\muk}^2\bigr)
\le \exp\!\left(-\frac{\norm{\mu_j-\muk}^4}{2\cdot4\sigma^2\norm{\mu_j-\muk}^2}\right)
= \exp\!\left(-\frac{\norm{\mu_j-\muk}^2}{8\sigma^2}\right),
\]
establishing the claim.  $\square$

Similarly, points from state $k$ are \emph{not} assigned to $V_k(\mu)$
only if they are closer to some $\mu_j$, $j\neq k$.  By the union bound
over $j\neq k$,
\[
P\bigl(\mu_k+\varepsilon\notin V_k(\mu)\mid S=k\bigr)
\;\le\; K\cdot\exp\!\left(-\frac{\Delta_{\min}^2}{8\sigma^2}\right).
\]

\noindent\textbf{Step 4 (Bounding $\norm{\ttheta_k-\muk}$).}
Decompose the conditional expectation:
\[
\ttheta_k = \E[\phi\mid\phi\in V_k(\mu)]
          = \frac{\E[\phi\cdot\ind\{\phi\in V_k(\mu)\}]}
                 {P(V_k(\mu))}.
\]
Write $\phi=\mu_S+\varepsilon$ and split by true state $S$:
\[
\E[\phi\cdot\ind\{\phi\in V_k(\mu)\}]
= \sum_{j=1}^K \pi_j\,
  \E\bigl[(\mu_j+\varepsilon)\cdot\ind\{\mu_j+\varepsilon\in V_k(\mu)\}\mid S=j\bigr].
\]
For the diagonal term ($j=k$),
\begin{align*}
\E[(\mu_k+\varepsilon)\cdot\ind\{\mu_k+\varepsilon\in V_k(\mu)\}\mid S=k]
&= \muk\cdot P(\mu_k+\varepsilon\in V_k(\mu)\mid S=k) \\
&\qquad + \E[\varepsilon\cdot\ind\{\mu_k+\varepsilon\in V_k(\mu)\}\mid S=k].
\end{align*}
For the off-diagonal terms ($j\neq k$), the indicator event has
probability $\le\exp(-\Delta_{\min}^2/(8\sigma^2))$.

The conditional expectation of the noise is bounded via H\"older's
inequality and the sub-Gaussian tail.  Let
$A_k=\{\mu_k+\varepsilon\notin V_k(\mu)\}$ be the event that a point from
state $k$ is assigned to a wrong Voronoi cell.  From Step~3,
$P(A_k\mid S=k)\le K\cdot\exp(-\Delta_{\min}^2/(8\sigma^2))$.  Since
$\E[\varepsilon]=0$ unconditionally,
\[
0=\E[\varepsilon\ind\{\mu_k+\varepsilon\in V_k(\mu)\}\mid S=k]
  +\E[\varepsilon\ind_{A_k}\mid S=k],
\]
hence
\[
\E[\varepsilon\mid\mu_k+\varepsilon\in V_k(\mu)]
=-\frac{\E[\varepsilon\ind_{A_k}\mid S=k]}
       {P(\mu_k+\varepsilon\in V_k(\mu)\mid S=k)}.
\]
By the Cauchy--Schwarz inequality,
\[
\norm{\E[\varepsilon\ind_{A_k}\mid S=k]}_2
\le\E[\norm{\varepsilon}_2\ind_{A_k}\mid S=k]
\le\E[\norm{\varepsilon}_2^2\mid S=k]^{1/2}\cdot P(A_k\mid S=k)^{1/2}.
\]
For the sub-Gaussian vector $\varepsilon$,
$\E[\norm{\varepsilon}_2^2]\le C_3 d_\phi\sigma^2$ (Vershynin~\cite{vershynin2018highdimensional},
Theorem~3.1.1), giving
\[
\norm{\E[\varepsilon\ind_{A_k}\mid S=k]}_2
\le\sqrt{C_3 K d_\phi}\,\sigma\cdot
  \exp\!\left(-\frac{\Delta_{\min}^2}{16\sigma^2}\right).
\]
Since $P(\mu_k+\varepsilon\in V_k(\mu)\mid S=k)
\ge 1-K\exp(-\Delta_{\min}^2/(8\sigma^2))\ge 1/2$ under strong separation,
\[
\norm{\E[\varepsilon\mid\mu_k+\varepsilon\in V_k(\mu)]}_2
\le 2\sqrt{C_3 K d_\phi}\,\sigma\cdot
   \exp\!\left(-\frac{\Delta_{\min}^2}{16\sigma^2}\right).
\]
(If a cruder bound with exponent $\Delta_{\min}^2/(8\sigma^2)$ is
desired, one may use the H\"older inequality with
$\norm{\varepsilon}_p$ and $p=\log(1/\delta)^{-1}$; the factor~$2$ in
the exponent is immaterial under \eqref{eq:S5:strongsep} with
$C_0$ large enough.)

Assembling terms,
\[
\norm{\ttheta_k-\muk}
\le \frac{C_3\,K\,(M+\sigma\sqrt{d_\phi})}{\pi_k}
    \cdot\exp\!\left(-\frac{\Delta_{\min}^2}{8\sigma^2}\right)
\le \epsilon_{\mathrm{pop}}.
\]

\noindent\textbf{Step 5 (Direct lower bound: global minimizer is near $\mu$).}
We now show that any $\theta$ far from $\mu$ incurs strictly larger
population risk than $W(\mu)$, so the global minimizer $\theta^*$ must
lie near $\mu$.  This avoids the circularity of assuming $\theta^*$ is
already in the contraction region.

Decompose the population risk as
\[
W(\theta)=\sum_{k=1}^K\pi_k\,
          \E\bigl[\norm{\mu_k+\varepsilon-\theta_{\pi(k)}}^2\cdot
                   \ind\{\mu_k+\varepsilon\in V_{\pi(k)}(\theta)\}\mid S=k\bigr]
          +\text{cross terms}.
\]
The cross terms (points from state $k$ assigned to a wrong center) each
have probability $\le\exp(-\Delta_{\min}^2/(8\sigma^2))$ by Step~3, so
their total contribution is $O(K(M^2+\sigma^2)\exp(-\Delta_{\min}^2/(8\sigma^2)))$.

Now suppose $\norm{\theta-\mu}_\infty\ge\Delta_{\min}/4$.  Then there
exists some $k$ such that $\norm{\theta_{\pi(k)}-\mu_k}\ge\Delta_{\min}/4$.
For points from state $k$ with mild noise
$\norm{\varepsilon}\le\Delta_{\min}/8$ (an event of probability
$1-2\exp(-c_0\Delta_{\min}^2/\sigma^2)$ by Lemma~\ref{lem:S5:sgnorm}),
the nearest center is $\theta_{\pi(k)}$ (since any competing center
$\theta_{\pi(j)}$ satisfies
$\norm{\mu_k+\varepsilon-\theta_{\pi(j)}}\ge
 \norm{\mu_k-\mu_j}-\norm{\theta_{\pi(j)}-\mu_j}-\norm{\varepsilon}
 \ge\Delta_{\min}-\Delta_{\min}/4-\Delta_{\min}/8>5\Delta_{\min}/8
 >\norm{\mu_k+\varepsilon-\mu_k}+\Delta_{\min}/4
 \ge\norm{\mu_k+\varepsilon-\theta_{\pi(k)}}$ by the separation
$\norm{\theta_{\pi(k)}-\mu_k}\ge\Delta_{\min}/4$ guarantees it is not
the closest).  Hence,
\[
W(\theta)-W(\mu)
\;\ge\; \pi_k\cdot\Bigl(
        \norm{\theta_{\pi(k)}-\mu_k}^2
        -\underbrace{\norm{\mu_k-\mu_k}^2}_{=0}
        \Bigr)\cdot\bigl(1-o(1)\bigr)
        \;-\; o(K(M^2+\sigma^2))
\;\ge\; \pi_{\min}\cdot\frac{\Delta_{\min}^2}{16}
        \;-\; o(1).
\]

Under \eqref{eq:S5:strongsep} with $C_0$ large enough, the gap
$\pi_{\min}\Delta_{\min}^2/16$ dominates the $o(1)$ correction,
so $\norm{\theta-\mu}_\infty\ge\Delta_{\min}/4$ implies
$W(\theta)>W(\mu)$.  Since $\theta^*$ is a global minimizer of $W$,
we must have $\norm{\theta^*-\mu}_\infty<\Delta_{\min}/4$.

To refine this to the sharper $\Delta_{\min}/8$ bound with the
exponentially small $\epsilon_{\mathrm{pop}}$, we combine the
self-consistency equation (1) with the fact that $\theta^*$ is already
known to be within $\Delta_{\min}/4$ of $\mu$.  For any such $\theta$,
the Voronoi partition is stable: points from state $k$ are assigned to
$\theta_{\pi(k)}$ with probability $1-O(\exp(-\Delta_{\min}^2/(8\sigma^2)))$,
and the self-consistency condition $\theta^*=T(\theta^*)$ implies
(Pollard~\cite{pollard1981strong}, Lemma~1):
\[
\norm{\theta^*_{\pi(k)}-\mu_k}
\;\le\; \frac{C_2\,K\,(M+\sigma\sqrt{d_\phi})}{\pi_k}
        \cdot\exp\!\left(-\frac{\Delta_{\min}^2}{8\sigma^2}\right)
\;=\; \epsilon_{\mathrm{pop}}.
\]
Under \eqref{eq:S5:strongsep} with $C_0$ sufficiently large,
$\epsilon_{\mathrm{pop}}<\Delta_{\min}/8$, because the exponential
$\exp(-\Delta_{\min}^2/(8\sigma^2))$ decays super-polynomially while
the pre-factor is polynomial.  The explicit condition is
\[
\frac{\Delta_{\min}^2}{8\sigma^2}\;\ge\;
\log\!\left(\frac{16C_2\,K\,(M+\sigma\sqrt{d_\phi})}
                 {\pi_{\min}\Delta_{\min}}\right),
\]
which is guaranteed by $C_0$ large enough.  \hfill$\square$
\end{proof}

\begin{remark}
The proof uses the self-consistency equation instead of directly
computing $W(\mu)-W(\mu^*)$, thereby avoiding the reversed-inequality
issue.  The bias $\epsilon_{\mathrm{pop}}$ is exponentially small in
$\Delta_{\min}^2/\sigma^2$ and is much smaller than the $\Delta_{\min}/8$
threshold needed for subsequent lemmas.
\end{remark}

% ---------------------------------------------------------------------------
\subsection{Lemma~2: Exponential Convergence of the Empirical Minimizer}
\label{sec:S5:lem2}
% ---------------------------------------------------------------------------

\begin{lemma}[Exponential convergence of empirical minimizer]
\label{lem:S5:emp}
Let $\theta^*$ be the unique population minimizer from
Lemma~\ref{lem:S5:pop} and let $\htheta_n$ be the global empirical
minimizer of $W_n$.  For any $t>0$ satisfying
$t\ge C_0\sqrt{K d_\phi/n}$, there exist constants
$c_2,C_4>0$ depending on the problem parameters such that
\begin{equation}
P\bigl(\norm{\htheta_n-\theta^*}\ge t\bigr)
\;\le\; C_4\cdot\exp\!\left(-c_2\cdot\frac{n\,t^2}{\sigma^2 d_\phi}\right).
\label{eq:S5:lem2_bound}
\end{equation}
The constants are $c_2=\lambda^2/(128 C_L^2)$ with
$\lambda=\pi_{\min}/2$ (the strong convexity parameter) and
$C_L=2(M+\sigma\sqrt{d_\phi})$ (the Lipschitz constant bound), and
$C_4=2$ (from the peeling sum).  In particular, for $t=\Delta_{\min}/8$
and $n\ge n_0$,
\[
P\!\left(\norm{\htheta_n-\theta^*}\ge\frac{\Delta_{\min}}{8}\right)
\;\le\; 2\cdot\exp\!\left(
          -c_2\cdot\frac{n\,\Delta_{\min}^2}{64\,\sigma^2 d_\phi}\right).
\]
\end{lemma}

\begin{proof}[Proof of Lemma~\ref{lem:S5:emp}]
The proof uses a localized argmin argument with peeling.  It has three
parts: (i) a quadratic lower bound for $W$ near $\theta^*$,
(ii~and~iii) control of the empirical process via covering numbers and
concentration, and (iv) a peeling argument to convert the empirical
process bound into a bound on $\norm{\htheta_n-\theta^*}$.

\noindent\textbf{Step 1 (Quadratic lower bound near $\theta^*$).}
Define $r_0=\Delta_{\min}/4$.  By Lemma~\ref{lem:S5:pop},
$\norm{\theta^*-\mu}<\Delta_{\min}/8$, so the set
$\{\theta:\norm{\theta-\theta^*}\le r_0\}$ is contained within
$3\Delta_{\min}/8$ of the true centers.  Within this region, strong
separation implies that for a point $\phi=\mu_k+\varepsilon$ with
$\norm{\varepsilon}<\Delta_{\min}/8$, the nearest center in $\theta$ is
uniquely $\theta_{\pi(k)}$ (the one matched to true state $k$).  Hence
the misassignment probability is exponentially small
(cf.\ Pollard~\cite{pollard1981strong}, Lemma~A, who establishes a
quadratic lower bound for $k$-means under well-separated mixtures).
A self-contained proof sketch under our strong separation condition is
in Lemma~\ref{lem:S5:quad} (Appendix).  The result is:
\[
W(\theta)-W(\theta^*)\;\ge\;\lambda\cdot\norm{\theta-\theta^*}^2,
\qquad \forall\;\norm{\theta-\theta^*}\le r_0,
\tag{2}
\]
where $\lambda=\pi_{\min}/2$.

\noindent\textbf{Step 2 (The argmin inequality).}
Let $\htheta_n$ minimize $W_n$.  From $W_n(\htheta_n)\le W_n(\theta^*)$,
\begin{align*}
0 &\ge W_n(\htheta_n)-W_n(\theta^*) \\
  &= \bigl(W(\htheta_n)-W(\theta^*)\bigr)
     + \bigl((P_n-P)(f_{\htheta_n}-f_{\theta^*})\bigr) \\
  &\ge \lambda\norm{\htheta_n-\theta^*}^2
     - \bigl|(P_n-P)(f_{\htheta_n}-f_{\theta^*})\bigr|,
\end{align*}
where $f_\theta(x)=\min_k\norm{\phi(x)-\theta_k}^2$.  Therefore,
\begin{equation}
\lambda\norm{\htheta_n-\theta^*}^2
\;\le\; \bigl|(P_n-P)(f_{\htheta_n}-f_{\theta^*})\bigr|.
\tag{3}
\end{equation}
This holds whenever $\norm{\htheta_n-\theta^*}\le r_0$; the case
$\norm{\htheta_n-\theta^*}>r_0$ is handled separately.

\noindent\textbf{Step 3 (The localized empirical process).}
Define the localized supremum
\[
\psi_n(r)=\sup_{\norm{\theta-\theta^*}\le r}
          \bigl|(P_n-P)(f_\theta-f_{\theta^*})\bigr|.
\]
From (3), if $\norm{\htheta_n-\theta^*}\le r_0$, then letting
$r=\norm{\htheta_n-\theta^*}$ gives
\begin{equation}
\lambda r^2 \;\le\; \psi_n(r).
\tag{4}
\end{equation}

\noindent\textbf{Step 4 (Expectation of $\psi_n(r)$).}
For a fixed radius $r$, consider the function class
\[
\mathcal{G}_r=\{\,g_\theta(x)=f_\theta(x)-f_{\theta^*}(x):
               \norm{\theta-\theta^*}\le r\,\}.
\]
Each $g_\theta$ satisfies $|g_\theta(x)|\le L(x)\,r$ where
$L(x)=2(\norm{\phi(x)}_2+M)$ (by the Lipschitz property,
Lemma~S5.2), and $\E[g_\theta(X)^2]\le C_L^2 r^2$ where
$C_L^2=8(\max_j\norm{\mu_j}^2+\sigma^2 d_\phi+M^2)$.

The covering number of $\Theta_r=\{\theta:\norm{\theta-\theta^*}\le r\}$
under $\norm{\cdot}_\infty$ is bounded (Lemma~S5.3):
\[
\log\mathcal{N}(\epsilon,\Theta_r,\norm{\cdot}_\infty)
\;\le\; K d_\phi\log\!\left(1+\frac{4r}{\epsilon}\right).
\]
By Dudley's entropy integral, the expected supremum satisfies
\begin{equation}
\E[\psi_n(r)]
\;\le\; C_5\cdot\sqrt{\frac{K d_\phi}{n}}\cdot C_L\cdot r,
\tag{5}
\end{equation}
where $C_5$ is a universal constant.

\noindent\textbf{Step 5 (Concentration of $\psi_n(r)$).}
For the function class $\mathcal{G}_r$ with $|g|\le C_L r$ and
$\E[g^2]\le C_L^2 r^2$, Talagrand's concentration inequality gives
\begin{equation}
P\bigl(\psi_n(r)\ge\E[\psi_n(r)]+u\bigr)
\;\le\; \exp\!\left(-\frac{c_6\,n\,u^2}{C_L^2\,r^2}\right),
\tag{6}
\end{equation}
for any $u\ge0$, where $c_6$ is a universal constant.

\noindent\textbf{Step 6 (Peeling argument).}
Fix $t$ such that $t\le r_0$ and
$t\ge\frac{8C_5 C_L}{\lambda}\sqrt{\frac{K d_\phi}{n}}$.
Let $J=\lceil\log_2(r_0/t)\rceil$ and define dyadic intervals
$r_j=2^j t$ for $j=0,1,\dots,J$.

From (4): if $\norm{\htheta_n-\theta^*}\ge t$, then
$\lambda\norm{\htheta_n-\theta^*}^2\le\psi_n(\norm{\htheta_n-\theta^*})$.
For $\norm{\htheta_n-\theta^*}$ falling in $[r_{j-1},r_j)$ for some $j$,
\[
\lambda r_{j-1}^2 \;\le\; \lambda\norm{\htheta_n-\theta^*}^2
\;\le\; \psi_n(\norm{\htheta_n-\theta^*}) \;\le\; \psi_n(r_j).
\]
Therefore,
\begin{equation}
P(\norm{\htheta_n-\theta^*}\ge t)
\;\le\; \sum_{j=0}^J P\bigl(\psi_n(r_j)\ge\lambda (r_{j-1})^2\bigr),
\tag{7}
\end{equation}
where $r_{-1}:=t/2$ for the $j=0$ term.

For each $j$, bound $P(\psi_n(r_j)\ge\lambda r_{j-1}^2)$.
Since $r_j=2r_{j-1}$, $\lambda r_{j-1}^2=\lambda r_j^2/4$.
From (5): $\E[\psi_n(r_j)]\le C_5 C_L\sqrt{K d_\phi/n}\,r_j$.
Set $u_j=\lambda r_j^2/4-\E[\psi_n(r_j)]$.  Our condition on $t$ ensures
\[
\E[\psi_n(r_j)]\le C_5 C_L\sqrt{K d_\phi/n}\,r_j
\le\frac{\lambda r_j^2}{8},
\]
since $r_j\ge t\ge\frac{8C_5 C_L}{\lambda}\sqrt{\frac{K d_\phi}{n}}$.
Hence $u_j\ge\lambda r_j^2/8$.

Applying (6) with $u=u_j$:
\begin{equation}
P\!\left(\psi_n(r_j)\ge\frac{\lambda r_j^2}{4}\right)
\;\le\; \exp\!\left(
          -\frac{c_6\,n\,(\lambda r_j^2/8)^2}{C_L^2\,r_j^2}\right)
= \exp\!\left(-\frac{c_6\,\lambda^2\,n\,r_j^2}{64\,C_L^2}\right).
\tag{8}
\end{equation}

Now $r_j^2=4^j t^2$ (for $j\ge0$) and $r_0=t$.  The sum in (7) becomes
\[
P(\norm{\htheta_n-\theta^*}\ge t)
\;\le\; \sum_{j=0}^J
       \exp\!\left(-\frac{c_6\,\lambda^2\,n\cdot4^j\,t^2}{64\,C_L^2}\right).
\]
The sum is dominated by the $j=0$ term because subsequent terms decay
geometrically:
\[
\sum_{j=0}^\infty \exp\!\left(-\frac{c_6\,\lambda^2\,n\cdot4^j\,t^2}
                                   {64\,C_L^2}\right)
\;\le\; 2\cdot\exp\!\left(-\frac{c_6\,\lambda^2\,n\,t^2}{64\,C_L^2}\right),
\]
for all $n$ and $t$ satisfying the threshold condition.  Therefore,
\[
P(\norm{\htheta_n-\theta^*}\ge t)
\;\le\; 2\cdot\exp\!\left(-\frac{c_2\,n\,t^2}{\sigma^2 d_\phi}\right),
\]
where $c_2=c_6\lambda^2/(128 C_L^2)$, using
$C_L^2\le C_7\sigma^2 d_\phi$ (the dominant term when
$\sigma^2 d_\phi\gg\max\norm{\mu}^2$, or a constant otherwise) to
absorb numerical constants.  This establishes \eqref{eq:S5:lem2_bound}.
\hfill$\square$
\end{proof}

\begin{remark}
The exponent $-\frac{c_2 n t^2}{\sigma^2 d_\phi}$ is explicitly negative
for all $n>0$, $t>0$.  There is no polynomial pre-factor in $n$ (the
peeling sum gives a factor of $2$, independent of $n$), guaranteeing
that the exponential dominates for large $n$.  The threshold
$t\ge C_0\sqrt{K d_\phi/n}$ is the statistical resolution limit.
\end{remark}

% ---------------------------------------------------------------------------
\subsection{Lemma~3: Deterministic Partition Recovery}
\label{sec:S5:lem3}
% ---------------------------------------------------------------------------

\begin{lemma}[Center proximity implies correct partition]
\label{lem:S5:part}
Let $\mu=\mutrue$ be the true centers with separation $\Delta_{\min}>0$.
Let $\theta=\{\theta_1,\dots,\theta_K\}$ be any set of estimated centers
satisfying, for some permutation $\pi$,
\[
\norm{\theta_j-\mu_{\pi(j)}}_2 \;\le\; \frac{\Delta_{\min}}{8}
\qquad \text{for all } j.
\]
Then for any point $\phi=\muk+\varepsilon$ with
$\norm{\varepsilon}_2<3\Delta_{\min}/8$,
\[
\argmin_{j}\norm{\phi-\theta_j}_2 \;=\; \pi^{-1}(k),
\]
i.e.\ the point is correctly classified by nearest-neighbour assignment to $\theta$.
\end{lemma}

\begin{proof}[Proof of Lemma~\ref{lem:S5:part}]
Fix a true state $k$.  Let $j_k=\pi^{-1}(k)$ be the estimated center
matched to $\muk$.  For any other estimated center $j'\neq j_k$, let
$k'=\pi(j')\neq k$ be its matched true center.

By the triangle inequality, the distance from $\phi$ to the correct
estimated center is
\[
\norm{\phi-\theta_{j_k}}
\;\le\; \norm{\muk-\theta_{j_k}} + \norm{\varepsilon}
\;\le\; \frac{\Delta_{\min}}{8} + \norm{\varepsilon}.
\]

The distance to a wrong estimated center is
\begin{align*}
\norm{\phi-\theta_{j'}}
&\ge \norm{\muk-\mu_{k'}} - \norm{\mu_{k'}-\theta_{j'}} - \norm{\varepsilon} \\
&\ge \Delta_{\min} - \frac{\Delta_{\min}}{8} - \norm{\varepsilon}
= \frac{7\Delta_{\min}}{8} - \norm{\varepsilon}.
\end{align*}

For $\norm{\varepsilon}<3\Delta_{\min}/8$,
\begin{align*}
\norm{\phi-\theta_{j_k}}
&\le \frac{\Delta_{\min}}{8} + \frac{3\Delta_{\min}}{8}
= \frac{\Delta_{\min}}{2}, \\[4pt]
\norm{\phi-\theta_{j'}}
&\ge \frac{7\Delta_{\min}}{8} - \frac{3\Delta_{\min}}{8}
= \frac{\Delta_{\min}}{2}.
\end{align*}

Therefore $\norm{\phi-\theta_{j_k}}\le\Delta_{\min}/2\le
\norm{\phi-\theta_{j'}}$, with strict inequality when the noise bound
is strict or $\norm{\muk-\theta_{j_k}}<\Delta_{\min}/8$.  The nearest
estimated center to $\phi$ is $\theta_{j_k}=\theta_{\pi^{-1}(k)}$, which
is the center matched to the true state $k$.  \hfill$\square$
\end{proof}

\begin{corollary}[Misclassification bound]
\label{cor:S5:misc}
Under the same conditions, the overall misclassification rate satisfies
\[
\frac1n\sum_{i=1}^n\ind\{\hat s(x_i)\neq s(x_i)\}
\;\le\; \ind\!\left\{\max_j\norm{\hat\theta_j-\mu_{\pi(j)}}
                \ge\frac{\Delta_{\min}}{4}\right\}
\;+\; \frac1n\sum_{i=1}^n\ind\{\norm{\varepsilon_i}\ge3\Delta_{\min}/8\}.
\]
\end{corollary}

\begin{proof}
The first term captures the event that the center proximity condition of
Lemma~\ref{lem:S5:part} fails.  The second term captures the irreducible
noise: even with perfect centers, a point with
$\norm{\varepsilon}\ge3\Delta_{\min}/8$ could be misclassified. \hfill$\square$
\end{proof}

\begin{remark}
The inequalities are verified:
$\frac18+\frac38=\frac12$ and $\frac78-\frac38=\frac12$, giving
$\norm{\phi-\theta_{j_k}}\le\Delta_{\min}/2\le\norm{\phi-\theta_{j'}}$
with the correct direction.  The proof is purely deterministic, and the
irreducible error is captured explicitly by the
$\norm{\varepsilon}\ge3\Delta_{\min}/8$ event.
\end{remark}

% ---------------------------------------------------------------------------
\subsection{Lemma~4: Lloyd's Algorithm Under Strong Separation}
\label{sec:S5:lem4}
% ---------------------------------------------------------------------------

The $k$-means problem is NP-hard in the worst case
\cite{aloise2009np,dasgupta2008hardness}.  The global empirical minimizer
$\htheta_n$ defined in Definition~\ref{def:S5:kmeans} is a computational
oracle.  In practice, Lloyd's algorithm (alternating assignment and update)
finds a local minimum.  Lemma~\ref{lem:S5:lloyd} shows that under the
strong separation condition \eqref{eq:S5:strongsep}, this computational gap
closes: the landscape is benign and Lloyd's with random restarts finds the
global minimizer with high probability.

\begin{lemma}[Lloyd's with random restarts under strong separation]
\label{lem:S5:lloyd}
Under the conditions of Theorem~\ref{thm:S5:main} (fixed $K$, strong
separation \eqref{eq:S5:strongsep}, $n$ sufficiently large), Lloyd's
algorithm initialized with $R=C_R\log n$ independent random initializations
(e.g.\ $k$-means++ or uniform subsampling of data points) returns a
solution $\ttheta_n$ satisfying
\[
P\!\left(\norm{\ttheta_n-\htheta_n^*}\ge\frac{\Delta_{\min}}{16}\right)
\;\le\; n^{-c}
\]
for any desired $c>0$ (by choosing $C_R$ sufficiently large).  Here
$\htheta_n^*$ denotes the global empirical $k$-means minimizer.
\end{lemma}

\begin{proof}[Proof of Lemma~\ref{lem:S5:lloyd}]

\noindent\textbf{Step 1 (Landscape structure under strong separation).}
From Lemma~\ref{lem:S5:pop}, $\norm{\theta^*-\mu}<\Delta_{\min}/8$.
From Lemma~\ref{lem:S5:emp}, for $n$ large enough,
$\norm{\htheta_n^*-\theta^*}<\Delta_{\min}/8$ with probability
$1-\exp(-c n)$.  Hence $\norm{\htheta_n^*-\mu}<\Delta_{\min}/4$ with
high probability.

Now consider $W_n$ restricted to the ball
$B=\{\theta:\norm{\theta-\mu}_\infty\le\Delta_{\min}/4\}$.  Within this
ball:
\begin{itemize}
\item The Voronoi partition is stable: moving any center by at most
$\Delta_{\min}/8$ does not change the assignment of any point with
$\norm{\varepsilon}<\Delta_{\min}/4$ (by Lemma~\ref{lem:S5:part}).
The fraction of points with $\norm{\varepsilon}\ge\Delta_{\min}/4$ is
at most $2\exp(-\Delta_{\min}^2/(32\sigma^2))$, which is
$O(n^{-C_0/2})$ under \eqref{eq:S5:strongsep}.
\item Consequently, $W_n$ is \textbf{strongly convex} within $B$
(structurally, it is a sum of quadratic functions, one per cluster,
with Hessian $\succeq (n_{\min}/n)\cdot I$).
\item The unique stationary point of $W_n$ within $B$ is the global
minimizer $\htheta_n^*$.
\end{itemize}

\noindent\textbf{Step 2 (Lloyd's as alternating minimization).}
Lloyd's algorithm iterates:
\begin{enumerate}
\item \textbf{Assignment}: For fixed centers $\theta^{(t)}$, assign each
point to the nearest center.
\item \textbf{Update}: Set
$\theta_j^{(t+1)}=\frac1{|C_j^{(t)}|}\sum_{i\in C_j^{(t)}}\phi(x_i)$,
where $C_j^{(t)}$ is the set of points assigned to center $j$.
\end{enumerate}
Within the ball $B$, the assignment step correctly clusters all but an
exponentially small fraction of points (by Lemma~\ref{lem:S5:part}).
The update step therefore computes the sample mean of nearly-clean
clusters, which is a contraction toward the true cluster means:
\[
\norm{\theta_j^{(t+1)}-\htheta_{n,j}^*}
\;\le\; \frac12\norm{\theta_j^{(t)}-\htheta_{n,j}^*}
\]
for all $j$, with probability at least
$1-\exp(-c n\Delta_{\min}^2/(K\sigma^2 d_\phi))$.  Thus Lloyd's converges
linearly to $\htheta_n^*$ from any initialization in $B$, achieving
accuracy $\Delta_{\min}/16$ in $O(\log n)$ iterations.

\noindent\textbf{Step 3 (Random initialization covers the good basin).}
We analyze the probability that a single random initialization lands in
the good basin $B=\{\theta:\norm{\theta-\mu}_\infty\le\Delta_{\min}/4\}$.

For uniform subsampling initialization, we select $K$ data points
independently (with replacement) and use them as initial centers.  Let
$E_1$ be the event that at least one point from each true cluster is
selected, and $E_2$ the event that all $K$ selected points are within
$\Delta_{\min}/8$ of their respective true cluster centers.  The
initialization lands in $B$ only if $E_1\cap E_2$ occurs, and then we
pair each initial center with the nearest true center.

\noindent\textit{Bounding $P(E_1)$.}
$P(E_1)\ge\prod_{k=1}^K (1-(1-\pi_k)^K)\ge\prod_{k=1}^K\pi_k
\ge\pi_{\min}^K$.  This is a fixed constant for fixed $K$, but
decreases exponentially in $K$.

\noindent\textit{Bounding $P(E_2\mid E_1)$.}
Conditional on $E_1$, each selected point is a sample from its own
cluster, and by the sub-Gaussian tail bound (Lemma~\ref{lem:S5:sgnorm}),
the probability that all $K$ selected points satisfy
$\norm{\phi-\mu_{s(\phi)}}<\Delta_{\min}/8$ is at least
$1-K\cdot 2\exp(-c_0\Delta_{\min}^2/\sigma^2)$.

\noindent\textit{Combined bound.}
A single random initialization lands in $B$ with probability
$p_0\ge\pi_{\min}^K\cdot(1-K\cdot 2\exp(-c_0\Delta_{\min}^2/\sigma^2))
\ge\pi_{\min}^K/2$ under strong separation.  While this is positive,
it is exponentially small in $K$; for large $K$ the required number of
restarts $R\approx p_0^{-1}\log n$ may become impractical.  We
acknowledge this gap: the present analysis applies only to the
\textbf{fixed-$K$} regime.  For growing $K$, one should use
$k$-means++ initialization~\cite{arthur2007kmeanspp} which provides an
$O(\log K)$ approximation guarantee and typically samples centers
proportional to squared distances, yielding a much larger basin coverage.
Alternatively, a spectral initialization~\cite{kumar2010clustering}
can guarantee landing in $B$ with probability $1-o(1)$ in one attempt
under strong separation, at the cost of an SVD computation.

\noindent\textbf{Step 4 (Amplification via multiple restarts).}
With $R$ independent restarts,
\[
P(\text{no restart lands in }B) \;\le\; (1-p_0)^R.
\]
Setting $R=C_R\log n$ with $C_R\ge(c+1)/|\log(1-p_0)|$ gives
\[
P(\text{initialization failure}) \;\le\; n^{-c}.
\]
Each successful initialization converges to within $\Delta_{\min}/16$ of
$\htheta_n^*$ in $O(\log n)$ Lloyd iterations.  The total runtime is
$O(R\cdot n\cdot K\cdot d_\phi\cdot\log n)=O(n\log^2 n)$, which is
polynomial.  The solution with the smallest $W_n$ among the $R$ runs is
returned. \hfill$\square$
\end{proof}

\begin{remark}[The NP-hard gap]
\label{rem:S5:npgap}
Lemma~\ref{lem:S5:lloyd} explicitly proves that under strong separation,
Lloyd's with random restarts finds the global empirical minimizer.
The proof exploits the benign landscape induced by well-separated
clusters, not the worst-case NP-hardness of $k$-means.

\textbf{If the strong separation condition does not hold}, the result
degrades.  The landscape may have spurious local minima, and Lloyd's may
converge to a suboptimal solution.  In that case, the theorem would need
to be weakened to ``there exists a polynomial-time algorithm that finds
a $\delta$-approximation of the global minimizer,'' following
\cite{kumar2010clustering,ostrovsky2013effectiveness}.  We
\textbf{honestly acknowledge this gap} and do not claim a guarantee for
the weak separation regime.

In practice, we recommend $k$-means++ initialization
\cite{arthur2007kmeanspp}, which provides $O(\log K)$ approximation in
expectation and has better practical coverage than uniform initialization.
Lemma~\ref{lem:S5:lloyd}'s conclusion holds for $k$-means++ as well,
with potentially better constants.
\end{remark}

% ---------------------------------------------------------------------------
\subsection{Proof of Main Theorem}
\label{sec:S5:main}
% ---------------------------------------------------------------------------

\begin{theorem}[Fixed-$K$ state discovery consistency]
\label{thm:S5:main}
Let $K$ be fixed.  Suppose:
\begin{enumerate}
\item \textbf{Generative model}: Assumption~\ref{ass:S5:gen} holds.
\item \textbf{Strong separation}: Assumption~\ref{ass:S5:sep} holds with
the constant $C_0$ large enough for Lemma~\ref{lem:S5:pop}.
\item \textbf{Asymptotics}: $n\to\infty$ and
$n_{\min}=\min_k n_k\to\infty$ (every state has infinitely many samples
in the limit).
\item \textbf{Algorithm}: Lloyd's $k$-means with
$R=C_R\log n$ independent random initializations, returning the solution
with the smallest $W_n$.
\end{enumerate}
Then the estimated partition $\hC$ satisfies
\begin{equation}
P\bigl(\hC\neq\cstar\text{ up to permutation}\bigr)
\;\le\; K\cdot\exp\!\left(
          -c_1\cdot\frac{n_{\min}\,\Delta_{\min}^2}{\sigma^2 d_\phi}\right)
        \;+\; o(1),
\label{eq:S5:main}
\end{equation}
where $c_1>0$ is a universal constant independent of
$K,n,\sigma,\Delta_{\min}$.  The $o(1)$ term captures the exponentially
small bias from the population minimizer and the initialization failure
probability from Lemma~\ref{lem:S5:lloyd}.
\end{theorem}

\begin{proof}[Proof of Theorem~\ref{thm:S5:main}]
We combine the four lemmas.

\noindent\textbf{Step 1 (Population bias is bounded).}
By Lemma~\ref{lem:S5:pop}, the unique population minimizer $\theta^*$
satisfies, for some permutation $\pi_1$,
\[
\norm{\theta^*_{\pi_1(k)}-\muk} \;\le\; \epsilon_{\mathrm{pop}}
\;<\; \frac{\Delta_{\min}}{8}.
\]

\noindent\textbf{Step 2 (Empirical minimizer convergence).}
By Lemma~\ref{lem:S5:emp} with $t=\Delta_{\min}/8$, for $n$ large
enough such that $\Delta_{\min}/8\ge C_0\sqrt{K d_\phi/n}$,
\[
P\!\left(\norm{\htheta_n-\theta^*}\ge\frac{\Delta_{\min}}{8}\right)
\;\le\; 2\cdot\exp\!\left(
          -\frac{c_2}{64}\cdot\frac{n\,\Delta_{\min}^2}{\sigma^2 d_\phi}\right).
\]
Combining with Lemma~\ref{lem:S5:pop} via the triangle inequality,
\[
\norm{\htheta_n-\mu}
\;\le\; \norm{\htheta_n-\theta^*} + \norm{\theta^*-\mu}
\;<\; \frac{\Delta_{\min}}{8} + \frac{\Delta_{\min}}{8}
= \frac{\Delta_{\min}}{4},
\]
with probability at least
$1-2\cdot\exp\!\left(-\frac{c_2}{64}\cdot\frac{n\,\Delta_{\min}^2}{\sigma^2 d_\phi}\right)$.

\noindent\textbf{Step 3 (Lloyd's finds the global minimizer).}
By Lemma~\ref{lem:S5:lloyd}, with $R=C_R\log n$ restarts,
\[
P\!\left(\norm{\ttheta_n-\htheta_n}\ge\frac{\Delta_{\min}}{16}\right)
\;\le\; n^{-c}
\]
for any desired $c>0$.  Since $\norm{\htheta_n-\mu}<\Delta_{\min}/4$
with high probability from Step~2, the Lloyd output satisfies
\[
\norm{\ttheta_n-\mu}
\;\le\; \norm{\ttheta_n-\htheta_n}
      + \norm{\htheta_n-\theta^*}
      + \norm{\theta^*-\mu}
\;<\; \frac{\Delta_{\min}}{16}
     + \frac{\Delta_{\min}}{8}
     + \frac{\Delta_{\min}}{8}
= \frac{\Delta_{\min}}{4},
\]
with probability at least
$1-2\exp\!\left(-c_2 n\Delta_{\min}^2/(64\sigma^2 d_\phi)\right)-n^{-c}$.

\noindent\textbf{Step 4 (Center proximity implies correct partition).}
By Lemma~\ref{lem:S5:part}, when
$\norm{\ttheta_j-\mu_{\pi(j)}}<\Delta_{\min}/4$ for all $j$, the induced
partition matches the true partition for all points with
$\norm{\varepsilon}<3\Delta_{\min}/8$.

The fraction of points with $\norm{\varepsilon}\ge3\Delta_{\min}/8$ is
bounded by the sub-Gaussian tail bound (Lemma~S5.4):
\[
P(\norm{\varepsilon}\ge3\Delta_{\min}/8)
\;\le\; 2\exp\!\left(-c_0\cdot\frac{\Delta_{\min}^2}{\sigma^2}\right).
\]
This irreducible error is independent of $n$ and is absorbed into the
$o(1)$ term.  It represents points that are inherently ambiguous due to
large noise---even perfect knowledge of the centers cannot classify them
correctly.

\noindent\textbf{Step 5 (Converting $n$ to $n_{\min}$).}
Since $K$ is fixed, $n\ge K\cdot n_{\min}$ (by the pigeonhole principle).
Therefore,
\[
\exp\!\left(-\frac{c_2}{64}\cdot\frac{n\,\Delta_{\min}^2}{\sigma^2 d_\phi}\right)
\;\le\; \exp\!\left(-\frac{c_2 K}{64}\cdot
              \frac{n_{\min}\,\Delta_{\min}^2}{\sigma^2 d_\phi}\right).
\]

\noindent\textbf{Step 6 (Final probability bound).}
Let $c_1 = c_2 K / 64$.  Collecting all error terms,
\begin{align*}
&P(\text{estimated partition}\neq\text{true partition}) \\
&\le P\!\left(\norm{\ttheta_n-\mu}\ge\frac{\Delta_{\min}}{4}\right)
   + P(\text{Lloyd's initialization fails}) \\
&\le 2\cdot\exp\!\left(-c_2\cdot\frac{n\,\Delta_{\min}^2}{64\,\sigma^2 d_\phi}\right)
   \;+\; n^{-c} \\
&\le 2\cdot\exp\!\left(-c_1\cdot\frac{n_{\min}\,\Delta_{\min}^2}
                                {\sigma^2 d_\phi}\right)
   \;+\; o(1).
\end{align*}
The $o(1)$ term absorbs:
\begin{itemize}
\item The exponentially small bias $\epsilon_{\mathrm{pop}}$ from
Lemma~\ref{lem:S5:pop} (which ensures
$\norm{\theta^*-\mu}<\Delta_{\min}/8$).
\item The initialization failure probability $n^{-c}$ from
Lemma~\ref{lem:S5:lloyd}.
\item The irreducible misclassification from large-noise samples
($\norm{\varepsilon}\ge3\Delta_{\min}/8$), which is
$O(\exp(-c\Delta_{\min}^2/\sigma^2))$ and vanishes as
$\Delta_{\min}/\sigma\to\infty$.
\item The factor $K$ in front of the exponential (from the union bound
over $K$ centers) can be absorbed into the exponential via a
logarithmic correction that is $o(1)$.
\end{itemize}
This completes the proof of Theorem~\ref{thm:S5:main}.  \hfill$\square$
\end{proof}

% ---------------------------------------------------------------------------
\subsection{Discussion: NP-Hard Gap and Practical Implications}
\label{sec:S5:disc}
% ---------------------------------------------------------------------------

The following remarks situate Theorem~\ref{thm:S5:main} within the
broader theory and address the computational-statistical gap.

\begin{remark}[The NP-hard gap (resolution of Review Issue~3)]
\label{rem:S5:gap}
The proof chain resolves the NP-hard gap within the strong separation
regime:
\begin{enumerate}
\item Lemma~\ref{lem:S5:emp} assumes access to the \textbf{global
empirical minimizer} $\htheta_n$, which is a computational oracle.
\item Lemma~\ref{lem:S5:lloyd} proves that \textbf{Lloyd's algorithm
with $R=O(\log n)$ random restarts finds $\htheta_n$} with probability
$1-o(1)$ under \eqref{eq:S5:strongsep}.
\end{enumerate}
The chain holds because $k$-means on well-separated mixtures is not a
hard instance---the landscape has a unique local minimum within the
basin of attraction, and Lloyd's contracts toward it.  The NP-hardness
of $k$-means applies to worst-case instances, not the well-separated
Gaussian-like mixtures considered here.

\textbf{Without strong separation}: the landscape may have spurious
local minima, and Lloyd's may converge to a suboptimal solution.  In
that case, only weaker guarantees are available---e.g.\ there exists a
polynomial-time algorithm that finds a $(1+\delta)$-approximation of
the global minimizer \cite{kumar2010clustering,ostrovsky2013effectiveness}.
The theorem makes no claim for the weak separation regime.
\end{remark}

\begin{remark}[Relation to main text bounds]
\label{rem:S5:relation}
Theorem~\ref{thm:S5:main} complements the main text's negative results:
\begin{itemize}
\item Theorem~1~\eqref{eq:main:F1} gives a lower bound on the F1 score
under the detection guarantee; Theorem~5 gives the convergence rate for
the underlying partition recovery.
\item Theorem~4'~\eqref{eq:main:minimax} establishes the minimax lower
bound; Theorem~5 shows it is achievable up to constant factors under
strong separation.
\item Theorem~2~\eqref{eq:main:weakF1} gives a weak-feature F1 bound that
cannot exceed a loss baseline; the phase transition between
Theorem~2 and Theorem~5 is governed by the signal-to-noise ratio
$\Delta_{\min}^2/(\sigma^2 d_\phi)$.
\end{itemize}
\end{remark}

\begin{remark}[Fixed $K$ and growing-$K$ regimes]
\label{rem:S5:fixedK}
The proof relies on $K$ being fixed.  This allows:
\begin{itemize}
\item The covering number $\log\mathcal{N}(\epsilon)\le
K d_\phi\log(1+2M/\epsilon)$, which is $O(\log(1/\epsilon))$ for fixed $K$.
\item The strong convexity parameter $\lambda=\pi_{\min}/2$ is fixed
(does not decay with $K$).
\item The random initialization probability $p_0\ge\pi_{\min}^K$ is a
fixed constant.
\end{itemize}
For $K\to\infty$ as $n\to\infty$, the covering number grows, $\pi_{\min}$
could be $O(1/K)$ making $\lambda=O(1/K)$, and the initialization
probability decays exponentially in $K$.  The growing-$K$ regime requires
a separate analysis; see \cite{pollard1981strong,lei2013statistical}.
\end{remark}

\begin{remark}[Practical sample size guide]
\label{rem:S5:samples}
From \eqref{eq:S5:main}, to achieve misclassification probability
$\le\delta$ due to center estimation, it suffices to have
\[
n_{\min}\;\ge\;
\frac{C_1\,\sigma^2 d_\phi}{\Delta_{\min}^2}
\cdot\log\!\left(\frac{K}{\delta}\right),
\]
for a universal constant $C_1>0$.  The following operational rule of
thumb illustrates the required per-state sample sizes:
\[
\begin{array}{c|c}
\Delta_{\min}/(\sigma\sqrt{d_\phi}) & n_{\min}\text{ for }P(\text{error})\le0.05 \\ \hline
0.5\ \text{(marginal)} & \ge400 \\
1.0\ \text{(moderate)} & \ge100 \\
2.0\ \text{(good)}     & \ge25  \\
3.0\ \text{(strong)}   & \ge12
\end{array}
\]
(Assuming $K\le10$, $d_\phi\le64$, and a refined constant $C_1\approx20$.)
\end{remark}

% ---------------------------------------------------------------------------
% Verification table
% ---------------------------------------------------------------------------
\subsection{Verification of Exponents and Inequality Directions}
\label{sec:S5:verify}

We explicitly verify that all exponents are negative and all inequality
directions are correct.

\medskip
\begin{center}
\begin{tabular}{lcc}
\toprule
\textbf{Step} & \textbf{Inequality} & \textbf{Exponent negative?} \\
\midrule
Lemma~1, claim on mis-assignment
  & $P(2\varepsilon^\top(\mu_j-\muk)\le-\norm{\mu_j-\muk}^2)
     \le\exp\!\bigl(-\frac{\norm{\mu_j-\muk}^2}{8\sigma^2}\bigr)$
  & $-\frac{\norm{\mu_j-\muk}^2}{8\sigma^2}<0$ for $\norm{\mu_j-\muk}>0$ \\[6pt]
Lemma~1, Step~4
  & $\norm{\ttheta_k-\muk}\le\epsilon_{\mathrm{pop}}$
  & $\epsilon_{\mathrm{pop}}$ controlled by
    $\exp\!\bigl(-\frac{\Delta_{\min}^2}{8\sigma^2}\bigr)<1$ \\[6pt]
Lemma~2, eq.\ (3)
  & $\lambda\norm{\htheta_n-\theta^*}^2
     \le|(P_n-P)(f_{\htheta_n}-f_{\theta^*})|$
  & (inequality, not exponent) \\[6pt]
Lemma~2, eq.\ (5)
  & $\E[\psi_n(r)]\le
     C_5\sqrt{\frac{K d_\phi}{n}}\,C_L r$
  & (inequality, not exponent) \\[6pt]
Lemma~2, eq.\ (8)
  & $P(\psi_n(r_j)\ge\frac{\lambda r_j^2}{4})
     \le\exp\!\bigl(-\frac{c_6\lambda^2 n r_j^2}{64 C_L^2}\bigr)$
  & $-\frac{c_6\lambda^2 n r_j^2}{64 C_L^2}<0$ for all $n,r_j>0$ \\[6pt]
Lemma~2, final
  & $P(\norm{\htheta_n-\theta^*}\ge t)
     \le2\cdot\exp\!\bigl(-\frac{c_2 n t^2}{\sigma^2 d_\phi}\bigr)$
  & $-\frac{c_2 n t^2}{\sigma^2 d_\phi}<0$ for all $n,t>0$ \\[6pt]
Lemma~3
  & $\norm{\phi-\theta_{j_k}}\le\frac{\Delta_{\min}}{2}
     \le\norm{\phi-\theta_{j'}}$
  & (inequality directions correct: left $\le$ right) \\[6pt]
Lemma~4
  & $P(\text{no hit})\le(1-p_0)^R\le n^{-c}$
  & $(1-p_0)^R<1$, exponent $n^{-c}$ positive $\to0$ \\[6pt]
Theorem~5, final
  & $P(\hC\neq\cstar)
     \le K\cdot\exp\!\bigl(-c_1\frac{n_{\min}\Delta_{\min}^2}
                                 {\sigma^2 d_\phi}\bigr)+o(1)$
  & $-\frac{c_1 n_{\min}\Delta_{\min}^2}{\sigma^2 d_\phi}<0$
    for all $n_{\min}>0$, $\Delta_{\min}>0$,
    $\sigma^2<\infty$, $d_\phi<\infty$ \\[6pt]
\bottomrule
\end{tabular}
\end{center}

The exponent $-\frac{c_1 n_{\min}\Delta_{\min}^2}{\sigma^2 d_\phi}$ is
\textbf{always negative} for $n_{\min}>0$, $\Delta_{\min}>0$,
$\sigma^2<\infty$, $d_\phi<\infty$.  There is no risk of a positive
exponent.  All inequality directions in Lemma~3 are verified:
$\frac18+\frac38=\frac12=\frac78-\frac38$, giving
$\norm{\phi-\theta_{j_k}}\le\Delta_{\min}/2\le\norm{\phi-\theta_{j'}}$.
The union bound in Theorem~\ref{thm:S5:main} propagates these correctly.

% ---------------------------------------------------------------------------
\subsection*{Appendix to Section~S5: Supporting Lemmas}
\label{sec:S5:app}
% ---------------------------------------------------------------------------

The following lemmas are used in the proofs above and are stated here
for completeness.

\begin{lemma}[Sub-Gaussian norm bound]
\label{lem:S5:sgnorm}
Let $\varepsilon\in\R^{d_\phi}$ be a zero-mean sub-Gaussian vector with
variance proxy $\sigma^2$.  For any $t>0$,
\[
P\!\left(\norm{\varepsilon}_2\ge C_1\sigma(\sqrt{d_\phi}+\sqrt{t})\right)
\;\le\; 2\exp(-t),
\]
where $C_1$ is an absolute constant.  For
$t=c\cdot\Delta_{\min}^2/\sigma^2$ with $\Delta_{\min}^2/\sigma^2$
large under \eqref{eq:S5:strongsep},
\[
P\!\left(\norm{\varepsilon}_2\ge\frac{\Delta_{\min}}{4}\right)
\;\le\; 2\exp\!\left(-c_0\frac{\Delta_{\min}^2}{\sigma^2}\right),
\]
for some $c_0>0$.  Moreover,
$P(\norm{\varepsilon}\ge3\Delta_{\min}/8)
\le 2\exp(-c_0'\Delta_{\min}^2/\sigma^2)$.
\end{lemma}

\begin{proof}
This follows from Theorem~3.1.1 of \cite{vershynin2018highdimensional}
for sub-Gaussian random vectors.  The second inequality follows by
setting $t=\sqrt{c}\,\Delta_{\min}/\sigma$ and noting that
$\sqrt{d_\phi}\le\Delta_{\min}/(C_1 C\sigma)$ under
\eqref{eq:S5:strongsep}.
\end{proof}

\begin{lemma}[Quadratic lower bound for $W$ near $\theta^*$]
\label{lem:S5:quad}
Under the conditions of Theorem~\ref{thm:S5:main}, for any $\theta$ with
$\norm{\theta-\theta^*}\le\Delta_{\min}/4$,
\[
W(\theta)-W(\theta^*)\;\ge\;\lambda\cdot\norm{\theta-\theta^*}^2,
\]
where $\lambda=\pi_{\min}/2$ and $\pi_{\min}=\min_k P(S=k)$.
\end{lemma}

\begin{proof}
Let $\theta$ satisfy $\norm{\theta-\theta^*}\le\Delta_{\min}/4$.  By
Lemma~\ref{lem:S5:pop}, $\norm{\theta^*-\mu}<\Delta_{\min}/8$, so
$\norm{\theta-\mu}<\Delta_{\min}/8+\Delta_{\min}/4=3\Delta_{\min}/8$.
For a point $\phi=\muk+\varepsilon$ with $\norm{\varepsilon}<\Delta_{\min}/8$,
Lemma~\ref{lem:S5:part} shows the nearest center in $\theta$ is the one
matched to true state~$k$.  Points with $\norm{\varepsilon}\ge\Delta_{\min}/8$
have probability at most $2\exp(-\Delta_{\min}^2/(128\sigma^2))$, which is
$o(\pi_{\min})$ under \eqref{eq:S5:strongsep}.

In the region where assignments are correct, the contribution to $W$ from
state $k$ is
\[
\E[\norm{\phi-\theta_{\pi(k)}}^2\mid S=k]\cdot\pi_k
= \pi_k\bigl(\norm{\muk-\theta_{\pi(k)}}^2+\E[\norm{\varepsilon}^2]\bigr),
\]
since the cross term vanishes because $\E[\varepsilon]=0$ and
$\varepsilon\perp S$.  The first term is minimized at
$\theta_{\pi(k)}=\muk$ with value $\pi_k\cdot\E[\norm{\varepsilon}^2]$.
Summing over $k$ and accounting for exponentially small corrections from
misassigned points gives
$W(\theta)-W(\theta^*)\ge\pi_{\min}\norm{\theta-\mu}^2
 - O(\exp(-\Delta_{\min}^2/(128\sigma^2)))$.
Since $\norm{\theta^*-\mu}$ is exponentially small,
$\norm{\theta-\theta^*}\approx\norm{\theta-\mu}$, and the result
follows with the factor $1/2$ absorbing the negligible corrections.
\end{proof}

\begin{lemma}[Lipschitz property of $k$-means loss]
\label{lem:S5:lipschitz}
For any fixed $x$, the function
$\theta\mapsto f_\theta(x)=\min_k\norm{\phi(x)-\theta_k}_2^2$ is
$L(x)$-Lipschitz in $\theta$ with respect to $\norm{\cdot}_\infty$,
where $L(x)=2(\norm{\phi(x)}_2+M)$.
\end{lemma}

\begin{proof}
For two center sets $\theta,\theta'$ with
$\max_k\norm{\theta_k-\theta'_k}\le\delta$,
\begin{align*}
|f_\theta(x)-f_{\theta'}(x)|
&\le \max_k\bigl|\norm{\phi(x)-\theta_k}^2
            -\norm{\phi(x)-\theta'_k}^2\bigr| \\
&\le \max_k\bigl(\norm{\phi(x)-\theta_k}+\norm{\phi(x)-\theta'_k}\bigr)
    \cdot\norm{\theta_k-\theta'_k} \\
&\le 2(\norm{\phi(x)}_2+M)\,\delta,
\end{align*}
using the triangle inequality and $\norm{\theta_k}\le M$.
\end{proof}

\begin{lemma}[Metric entropy of the center class]
\label{lem:S5:covering}
Let $\Theta_K=\{\theta=\{\theta_1,\dots,\theta_K\}:\theta_k\in\R^{d_\phi},
\norm{\theta_k}_2\le M\}$ equipped with
$\norm{\theta-\theta'}_\infty=\max_k\norm{\theta_k-\theta'_k}_2$.
For any $\epsilon\in(0,M)$,
\[
\log\mathcal{N}(\epsilon,\Theta_K,\norm{\cdot}_\infty)
\;\le\; K d_\phi\log\!\left(1+\frac{2M}{\epsilon}\right).
\]
\end{lemma}

\begin{proof}
$\Theta_K$ is the Cartesian product of $K$ balls of radius $M$ in
$\R^{d_\phi}$.  The $\epsilon$-covering number of $B_M(0)\subset\R^{d_\phi}$
in $\ell_2$ is $(1+2M/\epsilon)^{d_\phi}$, and the $\ell_\infty$ metric
on the product is the max of per-coordinate $\ell_2$ distances, so the
covering number of the product is the product of per-coordinate covering
numbers.
\end{proof}

% ---------------------------------------------------------------------------
% References (Section S5 bibliography)
% ---------------------------------------------------------------------------
\begin{thebibliography}{20}

\bibitem{aloise2009np}
D.~Aloise, A.~Deshpande, P.~Hansen, and P.~Popat.
NP-hardness of Euclidean sum-of-squares clustering.
\emph{Machine Learning}, 75(2):245--248, 2009.

\bibitem{arthur2007kmeanspp}
D.~Arthur and S.~Vassilvitskii.
$k$-means++: The advantages of careful seeding.
In \emph{Proceedings of SODA}, pages 1027--1035, 2007.

\bibitem{dasgupta2008hardness}
S.~Dasgupta.
The hardness of $k$-means clustering.
Technical Report, UCSD, 2008.

\bibitem{kumar2010clustering}
A.~Kumar and R.~Kannan.
Clustering with spectral norm and the $k$-means algorithm.
In \emph{Proceedings of FOCS}, pages 299--308, 2010.

\bibitem{lei2013statistical}
J.~Lei, A.~Rinaldo, and L.~Wasserman.
A statistical analysis of clustering algorithms.
\emph{arXiv preprint arXiv:1306.6836}, 2013.

\bibitem{ostrovsky2013effectiveness}
R.~Ostrovsky, Y.~Rabani, L.~J.~Schulman, and C.~Swamy.
The effectiveness of Lloyd-type methods for the $k$-means problem.
\emph{Journal of the ACM}, 59(6):1--22, 2013.

\bibitem{pollard1981strong}
D.~Pollard.
Strong consistency of $k$-means clustering.
\emph{The Annals of Statistics}, 9(1):135--140, 1981.

\bibitem{vershynin2018highdimensional}
R.~Vershynin.
\emph{High-Dimensional Probability: An Introduction with Applications in
Data Science}. Cambridge University Press, 2018.

\end{thebibliography}
